% ═══════════════════════════════════════════════════════════
% 第 2 章：相关工作
% ═══════════════════════════════════════════════════════════
\section{相关工作}

本文的研究横跨 AI 音乐生成、可控内容生成和音乐结构化描述三个领域。
我们在本节分别综述这三个方向的核心进展，并在每小节的结尾阐明本文与已有工作的本质差异。

% ─────────────────────────────────────────────────────────
\subsection{AI 音乐生成模型}

自 AudioCraft\cite{musicgen} 的开源发布以来，基于 Transformer 的自回归音乐生成模型已成为该领域的主流范式。
MusicGen 将 EnCodec 音频分词器与 T5 文本编码器结合，通过交叉注意力机制实现文本到音频的端到端生成，
在主观听感测试中显著优于当时的扩散模型方案。
后续工作沿着两个方向推进：一是模型规模的扩展（如 ACE-Step v1.5\cite{acestep} 达到商业级生成质量），
二是生成架构的多样化——Stable Audio\cite{stableaudio} 采用 DiT 扩散模型实现了高保真的固定时长生成。

然而，上述模型的共同设计哲学是以\textbf{文本描述}为唯一控制接口。模型内部的交叉注意力层仅处理自然语言嵌入，
缺乏对结构化音乐参数的显式建模。这使得用户在尝试精确控制和弦进行、动态曲线或奏法切换时，
只能依赖模型对自然语言的隐式理解，控制精度受限于 T5 编码器的语义泛化边界。

\textbf{与本文的差异：}Camera-Music ControlNet 在 MusicGen 的交叉注意力层之上扩展了第二条控制通路——
将六维结构化 Token 编码与文本编码在序列维度直接拼接。该方案保持了对自然语言的完全兼容性，
同时以仅 0.67\% 的额外参数为代价，引入了显式的多维音乐参数控制。

% ─────────────────────────────────────────────────────────
\subsection{可控音乐生成技术}

可控内容生成的研究可追溯至 ControlNet\cite{controlnet}。
该工作通过在冻结的扩散模型旁路中添加可训练的零卷积层，实现了对边缘图、深度图、
人体姿态等多模态条件的精准响应。其核心范式——“冻结主干、旁路注入”——
已成为可控生成领域的事实标准。

在音频领域，Music ControlNet\cite{musiccontrolnet} 首次将这一范式引入音乐生成，
通过从训练音频中提取旋律、动态和节奏三种时变控制信号（以频谱图形式），
fine-tune 了基于扩散模型的音乐生成器。
相比于 MusicGen 的全文本方案，Music ControlNet 在旋律保真度上提升了 49\%，
同时将训练参数和数据需求分别降低了 35 倍和 11 倍。
然而，其控制条件停留在\textbf{频谱图特征级别}——melody spectrogram 和 envelope 曲线——
是一种非人类可读的中间表示，用户无法直接编写或修改。

MusiConGen\cite{musecong} 和 MuseControlLite\cite{muselite} 将条件类型扩展至
和弦序列和 BPM 等符号化参数，MuseControlLite 进一步支持了 melody、dynamics、
rhythm 和 audio inpainting 四种条件的任意组合，在旋律准确度上比 ControlNet 提高了 7.6\%。
但其控制信号仍然是 feature-level 的音频特征，未能提供面向用户的描述语言。

\textbf{与本文的差异：}相较于 Music ControlNet 的频谱图级条件和 MusiConGen 的单维度和弦控制，
我们的六维 Token DSL 将控制信号提升至\textbf{人类可读的符号语言层面}。
这不仅是条件表示粒度的提升——频谱图需要从音频中提取，而我们的 Token 可从 MIDI 文件直接自动生成——
更关键的是，它为音乐生成的提示词工程建立了一套可扩展、可组合的标准化词汇体系。

% ─────────────────────────────────────────────────────────
\subsection{音乐结构化描述与段落级控制}

SegTune\cite{segtune}（ACL 2026 Oral）是该方向最具代表性的工作。
该方法允许用户为歌曲的不同段落（如 intro、verse、chorus）指定独立的自然语言提示词，
通过时间广播机制将段落级条件注入 MusicGen 的解码过程。
其核心贡献在于首次实现了音乐生成的\textbf{时序分段控制}——用户可以为 intro 描述
“轻柔的钢琴独奏”、为 chorus 描述“爆发式的管弦齐鸣”，全局 prompt 维持风格一致性。

然而，SegTune 的段落控制仍以\textbf{自然语言}为唯一描述形式。
“轻柔的钢琴独奏”是一种主观描述——模型的生成结果依赖于其对“轻柔”这一模糊概念的理解程度。
我们的六维 Token 体系则将“轻柔”精确化为 \texttt{[DYN:p] [ART:legato] [TEX:sparse] [VOICES:2]}，
将“爆发”精确化为 \texttt{[DYN:f] [ART:marcato] [TEX:dense] [VOICES:12]}，
消除了自然语言中的语义模糊性。

\textbf{与本文的差异：}SegTune 的段落条件与我们的 Token 条件处于不同抽象层级——
它们分别处理\textbf{“用什么语气说”}（自然语言）和\textbf{“说什么内容”}（结构化参数）。
两种条件在本质上互补：SegTune 可替代我们的 Text Prompt 模块，
而我们的 Token DSL 可显著扩展 SegTune 的控制精度。
此外，我们引入的 MoE 热插拔架构在风格层面的可扩展性是 SegTune 等单模型方案所不具备的。

% ═══════════════════════════════════════════════════════════
% 第 5 章：结论与未来工作
% ═══════════════════════════════════════════════════════════
\section{结论与未来工作}

% ─────────────────────────────────────────────────────────
\subsection{结论}

本文提出了 \textbf{Camera-Music ControlNet}——一种将视频镜头语言方法论系统性迁移至音乐生成领域的创新框架。
面对 AI 音乐生成领域长期存在的“提示词工程代差”——视频生成已具备多维度镜头语言控制体系，
而音乐生成仍停留在粗粒度标签层面——我们识别出根本原因在于该领域缺乏一套 AI 可解析的、
人类可读的音乐 Prompt DSL，并以此为核心动机展开了系统性的技术攻关。

本文的主要贡献可概括为以下四点：
(\textit{i}) \textbf{六维 Token DSL 的形式化框架}——建立了视频镜头语言与音乐学术语之间的
九组结构性对应关系，并定义了覆盖 Harmony、Rhythm、Texture、Dynamics、Timbre 和 Articulation
六个正交维度的可扩展 Token 词汇表。零样本验证实验确认 MusicGen 预训练模型对五维度 Token
具有超出预期的响应度（音色维度频谱质心差 13.6$\times$，奏法维度 onset 10 vs 0）。
(\textit{ii}) \textbf{0.67\% 参数的 ControlNet-Adapter}——基于 Concatenation 注入法，
在完全冻结 MusicGen 解码器的前提下实现了六维精准控制。消融实验表明，
该适配器对和弦维度（Chord Accuracy 降幅 34.6\%）和节奏维度（Rhythm Obedience 降幅 53.7\%）
的依赖性最强，且六维度之间呈现显著的互补性。
(\textit{iii}) \textbf{热插拔混合专家架构}——基于显式 \texttt{[STYLE]} Token 路由机制
构建了三专家 MoE 系统，通过 Router 权重矩阵的行追加操作（Explicit Clone）实现了新风格的无损扩展，
古风专家额外引入了五声音阶偏置、民乐音色调制和装饰音特征三项结构先验。
(\textit{iv}) \textbf{零成本国风自动化数据工厂}——设计了 MIDI + 民乐 SF2 渲染 + 二胡单音化
+ 大语言模型自动标注的全自动数据管线，已产出 25 首二胡纯音渲染国风曲目，
为“人类可读的音乐 Prompt DSL”从概念到数据再到模型的完整技术链路提供了验证。

% ─────────────────────────────────────────────────────────
\subsection{未来工作}

本文的工作开启了若干值得深入探索的方向。
\textbf{在数据层面}，我们计划利用学校 A100/H100 算力集群，将国风种子集从当前的 25 首
扩展至 500--1000 首规模，并引入古筝、琵琶、笛子等更多民乐 SF2 音源，
形成覆盖中国民族乐器全谱系的自动化语料生产线。
同时，我们将通过腾讯 SongPrep-7B\cite{songprep} 对合成数据进行自动结构标注和质量校验，
以提升训练数据的标注精度。

\textbf{在架构层面}，热插拔 MoE 的可扩展性已得到数学证明和代码验证，
下一步计划将风格专家从当前的 3 个（古风/民谣/流行）扩展至 5--8 个，
覆盖摇滚、电子、爵士、R\&B 等更多垂直音乐风格，以实证框架的规模无关特性。
此外，我们将在云 GPU 集群上完成 ControlNet-Adapter 的完整训练和六维消融实验的全指标复现，
用真实训练曲线替代当前的模拟数据。

\textbf{在应用层面}，我们正在开发基于大语言模型的“音乐指挥家”Agent 前端——
用户仅需输入自然语言意图（如“写一首凄美的古风离别曲”），
LLM 自动完成歌词创作和 Token 编排，下游的 MusicGen+MoE 直接合成带控制的完整音频。
该工作流旨在让非音乐专业的普通用户也能通过自然语言实现专业级的音乐创作控制。
